\documentclass[11pt]{article}

% --- packages ---
\usepackage[utf8]{inputenc}
\usepackage[T1]{fontenc}
\usepackage{amsmath, amssymb, amsthm}
\usepackage{mathtools}
\usepackage{stmaryrd}
\usepackage[margin=1in]{geometry}
\usepackage{hyperref}
\usepackage{enumitem}

% --- theorem environments ---
\newtheorem{theorem}{Theorem}
\newtheorem{proposition}[theorem]{Proposition}
\newtheorem{lemma}[theorem]{Lemma}
\newtheorem{corollary}[theorem]{Corollary}
\newtheorem{conjecture}[theorem]{Conjecture}

\theoremstyle{definition}
\newtheorem{definition}[theorem]{Definition}
\newtheorem{question}[theorem]{Question}
\newtheorem{remark}[theorem]{Remark}

% --- macros ---
\newcommand{\R}{\mathbb{R}}
\newcommand{\N}{\mathbb{N}}
\newcommand{\E}{\mathbb{E}}
\newcommand{\Prob}{\mathbb{P}}
\newcommand{\A}{\mathcal{A}}
\newcommand{\F}{\mathcal{F}}
\newcommand{\T}{\mathcal{T}}
\newcommand{\Csem}{\mathcal{C}_{\mathrm{sem}}}
\newcommand{\Hom}{\mathrm{Hom}}
\newcommand{\Homp}{\mathrm{Hom}^{+}}
\newcommand{\RandomExt}{\mathrm{RandomExt}}
\newcommand{\avg}[2]{\llbracket #1 \rrbracket_{#2}}
\DeclareMathOperator{\supp}{supp}

\title{Equivalence of Forbidden-Subgraph Reasoning\\
       and Quotient Semantics in Flag Algebras}
\author{}
\date{}

\begin{document}
\maketitle

\begin{abstract}
We investigate two natural ways of expressing, within Razborov's flag-algebra
formalism, the statement that a flag-algebra element is asymptotically
non-negative under a forbidden-induced-subgraph constraint. The first works in the
quotient theory $\T' = \T + A$ where $A$ forbids a fixed induced subgraph; the
second works in the unconstrained theory $\T$ but conditions on root
homomorphisms that satisfy~$A$, and asks for almost-sure non-negativity along
ensembles of random homomorphisms (Razborov, Definition~10). We show that for
$\sigma = \emptyset$ these two formulations are equivalent, with the
$(\Leftarrow)$ direction relying essentially on the (induced) graph removal
lemma. We discuss what changes for general types $\sigma$ and identify the
mappings between homomorphism spaces that mediate the equivalence.
\end{abstract}

\section{Introduction and Setup}

Let $\T$ denote the theory $\T_{\mathrm{Graph}}$ of simple undirected graphs, and
let $A$ be a universal axiom that forbids a fixed induced subgraph $H$ (the
running example is $H = K_3$, in which case $A$ is the assertion
``$\phi(K_3) = 0$''). Write $\T' \coloneqq \T + A$. Following Razborov's
construction in~\S2.3.3 of~\cite{Razborov2007}, for each non-degenerate type
$\sigma$ in $\T'$ there is a canonical algebra homomorphism
\[
  \pi^{\sigma} \colon \A^{\sigma}[\T] \twoheadrightarrow \A^{\sigma}[\T']
\]
realising $\A^{\sigma}[\T']$ as a quotient of $\A^{\sigma}[\T]$. Concretely,
$\pi^{\sigma}$ sends a $\sigma$-flag $F$ to itself if its underlying model
satisfies $A$, and to $0$ otherwise.

We use the notation $\avg{\,\cdot\,}{\sigma}$ for Razborov's downward
averaging operator (\S2.2) and write $\langle \sigma \rangle \coloneqq
\avg{1_{\sigma}}{\sigma} \in \A^{\emptyset}[\T]$ for the unlabelled
counterpart of the identity flag.\footnote{We restrict $\sigma_0 = \emptyset$
throughout, so $\eta$ is the empty map and may be suppressed from the
notation.} For $\phi_0 \in \Homp(\A^{\emptyset}[\T], \R)$ with
$\phi_0(\langle \sigma \rangle) > 0$ we write $\RandomExt(\phi_0; \sigma)$ for
the unique probability measure on $\Homp(\A^{\sigma}[\T], \R)$ guaranteed by
Theorem~3.5 of~\cite{Razborov2007}; a sample from this measure is denoted
$\phi^{\sigma}$.

\medskip

The central question is whether two natural ways to assert ``non-negativity
under the constraint $A$'' coincide.

\begin{question}\label{q:main}
Let $f \in \A^{\sigma}[\T]$. Are the following equivalent?
\begin{enumerate}[label=\textup{(\arabic*)}, leftmargin=2.5em]
  \item\label{cond:quot}
        \emph{(Quotient semantics.)} For every
        $\phi \in \Homp(\A^{\sigma}[\T'], \R)$,
        \[ \phi(\pi^{\sigma}(f)) \;\geq\; 0. \]
  \item\label{cond:rand}
        \emph{(Ensemble semantics.)} For every
        $\phi_0 \in \Homp(\A^{\emptyset}[\T], \R)$ satisfying $\phi_0 \models A$
        (e.g.\ $\phi_0(K_3) = 0$) and $\phi_0(\langle \sigma \rangle) > 0$,
        \[
          \Prob_{\phi^{\sigma} \sim \RandomExt(\phi_0; \sigma)}
          \bigl[\phi^{\sigma}(f) \geq 0 \bigr] \;=\; 1.
        \]
\end{enumerate}
\end{question}

The motivation is practical: \ref{cond:rand} expresses ``$f \geq 0$ over
triangle-free graphs'' entirely within $\T$, never leaving the larger algebra
$\A^{\sigma}[\T]$. This is convenient when one wants to combine
forbidden-subgraph reasoning with other techniques that are naturally phrased
in $\T$.

\section{The Bridge: Mappings Between Homomorphism Spaces}\label{sec:bridge}

Two natural set-theoretic mappings underlie Question~\ref{q:main}.

\paragraph{(a) Lifting $\T'$-homomorphisms to $\T$-homomorphisms.}
Since $\pi^{\sigma}$ is surjective, every
$\phi \in \Homp(\A^{\sigma}[\T'], \R)$ pulls back to
\[
  \widetilde{\phi} \;\coloneqq\; \phi \circ \pi^{\sigma}
  \;\in\; \Homp(\A^{\sigma}[\T], \R),
\]
and $\widetilde{\phi}$ vanishes on every $\sigma$-flag whose underlying model
violates $A$. Conversely, any $\psi \in \Homp(\A^{\sigma}[\T], \R)$ with
$\psi(F) = 0$ for every $A$-violating $F$ descends to a unique homomorphism in
$\Homp(\A^{\sigma}[\T'], \R)$. Thus the lifting identifies
\[
  \Homp(\A^{\sigma}[\T'], \R)
  \;\cong\;
  \bigl\{\,\psi \in \Homp(\A^{\sigma}[\T], \R) : \psi \models A \,\bigr\}.
\]

\paragraph{(b) Forgetting the labelled vertices.}
For each $\sigma$ there is a ``forgetful'' map
\[
  \Phi^{\sigma} \colon
  \Homp(\A^{\sigma}[\T], \R) \to \Homp(\A^{\emptyset}[\T], \R),
  \qquad
  \Phi^{\sigma}(\psi)(g) \coloneqq \psi(\pi^{\sigma}(g))
  \text{ for } g \in \A^{\emptyset}[\T],
\]
where $\pi^{\sigma} \colon \A^{\emptyset}[\T] \to \A^{\sigma}[\T]$ is
Razborov's \emph{upward} operator from~\S2.3.1 (sending an unlabelled model
to the sum of all its labellings; see Theorem~2.8 of~\cite{Razborov2007} for
its key properties).\footnote{We overload the symbol $\pi^{\sigma}$ here: in
the rest of this note it denotes the quotient map
$\A^{\sigma}[\T] \to \A^{\sigma}[\T']$ of~\S2.3.3. Both maps are
algebra homomorphisms and the meaning is always clear from context.}
Crucially, if $\psi \models A$, then so does $\Phi^{\sigma}(\psi)$ on
$\A^{\emptyset}[\T]$, because $\pi^{\sigma}(F) = 0$ whenever the unlabelled
model $F$ violates $A$.

\paragraph{Compatibility with ensembles.}
The two maps fit together cleanly: given
$\phi_0 \in \Homp(\A^{\emptyset}[\T], \R)$ with $\phi_0(\langle \sigma \rangle)
> 0$, the ensemble $\RandomExt(\phi_0; \sigma)$ is supported on the fibre
$(\Phi^{\sigma})^{-1}(\phi_0)$, and is the canonical disintegration of $\phi_0$
along $\Phi^{\sigma}$ (Theorems~3.5 and~3.17 of~\cite{Razborov2007}).

\section{The Equivalence}

\begin{theorem}\label{thm:main}
Let $\sigma = \emptyset$. Then conditions~\ref{cond:quot}
and~\ref{cond:rand} of Question~\ref{q:main} are equivalent.
\end{theorem}

When $\sigma = \emptyset$ the ensemble is degenerate (a Dirac mass at
$\phi_0$ itself), and the equivalence reduces to the assertion
\[
  \pi^{\emptyset}(f) \in \Csem(\F^{\emptyset}[\T'])
  \iff
  \bigl(\phi_0 \models A \Rightarrow \phi_0(f) \geq 0\bigr),
\]
which is the natural statement of soundness/completeness of quotient
semantics. We will prove both directions and indicate where the argument
generalises to $\sigma \neq \emptyset$.

\subsection{Proof of \texorpdfstring{$\ref{cond:quot} \Rightarrow
\ref{cond:rand}$}{(1) implies (2)}}

This direction works for arbitrary $\sigma$.

\begin{proof}
Fix $\phi_0 \in \Homp(\A^{\emptyset}[\T], \R)$ with $\phi_0 \models A$ and
$\phi_0(\langle \sigma \rangle) > 0$, and let $\phi^{\sigma} \sim
\RandomExt(\phi_0; \sigma)$. We first show that $\phi^{\sigma}$ almost surely
satisfies $A$. Let $H$ be the forbidden induced subgraph, viewed as an
element of $\A^{\emptyset}[\T]$. By Theorem~2.8(a) of~\cite{Razborov2007}
(applied with $f = H$ and $g = 1_{\sigma}$),
\[
  \avg{\pi^{\sigma}(H)}{\sigma} \;=\; H \cdot \langle \sigma \rangle
  \quad \text{in } \A^{\emptyset}[\T].
\]
By the chain rule (Lemma~2.2 of~\cite{Razborov2007}), the product
$H \cdot \langle\sigma\rangle$ expands as a non-negative combination of
$\emptyset$-flags each of which contains $H$ as an induced subgraph; hence
$\phi_0(H) = 0$ forces $\phi_0(H \cdot \langle\sigma\rangle) = 0$. Therefore,
by Definition~10 of~\cite{Razborov2007},
\[
  \E\bigl[\phi^{\sigma}(\pi^{\sigma}(H))\bigr]
  \;=\;
  \frac{\phi_0\bigl(\avg{\pi^{\sigma}(H)}{\sigma}\bigr)}{\phi_0(\langle \sigma \rangle)}
  \;=\;
  \frac{\phi_0(H \cdot \langle \sigma \rangle)}{\phi_0(\langle\sigma\rangle)}
  \;=\; 0.
\]
Since $\phi^{\sigma}(\pi^{\sigma}(H)) \geq 0$ pointwise,
\[
  \Prob\bigl[\phi^{\sigma}(\pi^{\sigma}(H)) = 0\bigr] = 1.
\]
By Remark~5 in~\cite{Razborov2007} (the equality version of Theorem~3.18),
this propagates: $\phi^{\sigma}$ almost surely sends every $\sigma$-flag whose
underlying model contains $H$ to zero. Equivalently, $\phi^{\sigma}$ almost
surely satisfies $A$, so almost surely $\phi^{\sigma}$ descends to a member of
$\Homp(\A^{\sigma}[\T'], \R)$ via the identification of
Section~\ref{sec:bridge}\,(a).

Hypothesis~\ref{cond:quot} now applies to the descended homomorphism, giving
$\phi^{\sigma}(f) = \phi^{\sigma}(\pi^{\sigma}(f)) \geq 0$ almost surely (the
first equality holds whenever $\phi^{\sigma}$ vanishes on $\ker \pi^{\sigma}$,
i.e.\ almost surely).
\end{proof}

\subsection{Proof of \texorpdfstring{$\ref{cond:rand} \Rightarrow
\ref{cond:quot}$}{(2) implies (1)} for $\sigma = \emptyset$}

\begin{proof}
Let $\phi \in \Homp(\A^{\emptyset}[\T'], \R)$ be arbitrary; we must show
$\phi(\pi^{\emptyset}(f)) \geq 0$. By Theorem~3.3(b) of~\cite{Razborov2007}
applied to the theory $\T'$, there exists a convergent sequence
$\{G_n'\}_{n \in \N}$ of $H$-induced-free graphs with $|V(G_n')| \to \infty$
such that
\[
  \phi \;=\; \lim_{n \to \infty} p^{G_n'} \quad \text{in } \Homp(\A^{\emptyset}[\T'], \R).
\]
Viewing each $G_n'$ as a model of $\T$ (i.e.\ forgetting the axiom $A$) yields
a sequence in $\T$; let $\phi_0 \in \Homp(\A^{\emptyset}[\T], \R)$ be its
limit. By construction,
\[
  \phi_0(F) \;=\; \lim_{n \to \infty} p(F, G_n') \;=\; \phi(\pi^{\emptyset}(F))
  \quad \text{for every } F \in \F^{\emptyset}[\T],
\]
so $\phi_0$ satisfies $A$ (since each $G_n'$ is $H$-induced-free). For
$\sigma = \emptyset$ the ensemble $\RandomExt(\phi_0; \emptyset)$ is a Dirac
mass at $\phi_0$ itself, so~\ref{cond:rand} applied to this $\phi_0$ gives
$\phi_0(f) \geq 0$. Hence
\[
  \phi(\pi^{\emptyset}(f)) \;=\; \phi_0(f) \;\geq\; 0. \qedhere
\]
\end{proof}

\begin{remark}
The proof above only uses the existence of \emph{some} sequence of
$H$-induced-free graphs converging to $\phi$. It does \emph{not} require the
graph removal lemma. The removal lemma enters only if one wishes to start
from an arbitrary sequence in $\T$ with vanishing $H$-density and convert it
to an $H$-induced-free sequence; see Section~\ref{sec:rl}.
\end{remark}

\section{The Role of the Removal Lemma}\label{sec:rl}

There is an alternative, more constructive proof of $\ref{cond:rand}
\Rightarrow \ref{cond:quot}$ that proceeds directly inside $\T$ and never
appeals to Theorem~3.3(b) for $\T'$. We sketch it because it isolates the
exact role of the removal lemma and naturally generalises to other forbidden
configurations.

Let $\phi \in \Homp(\A^{\emptyset}[\T'], \R)$. Pull back to
$\widetilde{\phi} = \phi \circ \pi^{\emptyset} \in \Homp(\A^{\emptyset}[\T],
\R)$; this homomorphism satisfies $A$. By Theorem~3.3(b) for $\T$, there is a
convergent sequence $\{G_n\}$ of $\T$-models (i.e.\ arbitrary graphs) with
$p^{G_n} \to \widetilde{\phi}$. In particular $p(H, G_n) \to
\widetilde{\phi}(H) = 0$.

\paragraph{Removal step.} The graph removal lemma for induced subgraphs of
Alon, Fischer, Krivelevich, and Szegedy~\cite{AFKS2000} yields graphs
$G_n'$ on the same vertex set as $G_n$ such that:
\begin{itemize}[leftmargin=2em]
  \item $G_n'$ is $H$-induced-free, and
  \item $|E(G_n) \mathbin{\triangle} E(G_n')| = o(n^2)$.
\end{itemize}

\paragraph{Stability of densities under $o(n^2)$-perturbations.}
A standard counting argument shows that for any fixed graph $F$ on $\ell$
vertices,
\[
  \bigl|\,p(F, G_n) - p(F, G_n')\,\bigr|
  \;\leq\;
  \binom{\ell}{2} \cdot
  \frac{|E(G_n) \mathbin{\triangle} E(G_n')|}{\binom{n}{2}}
  \;=\; o(1).
\]
Hence $\{G_n'\}$ is a convergent sequence with the same limit
$\widetilde{\phi}$, and it is $H$-induced-free.

\paragraph{Conclusion.} Apply~\ref{cond:rand} to $\widetilde{\phi}$ (which
satisfies $A$). Since $\sigma = \emptyset$, the ensemble is the Dirac mass at
$\widetilde{\phi}$ itself, and we conclude $\widetilde{\phi}(f) \geq 0$,
i.e.\ $\phi(\pi^{\emptyset}(f)) \geq 0$. \qed

\medskip

The role of the removal lemma is therefore not strictly necessary for the
equivalence per se---Theorem~3.3(b) applied to $\T'$ already suffices---but it
becomes essential if one wishes to prove the equivalence by staying inside
$\T$ throughout, or if one wishes to relate sequences in $\T$ with vanishing
$H$-density to sequences in $\T'$.

\begin{remark}[Generalisation to arbitrary hereditary properties]
The same argument works whenever $A$ is the conjunction of forbidding finitely
many (or even infinitely many) induced subgraphs, provided the corresponding
induced removal lemma holds. By the theorem of Alon and
Shapira~\cite{AlonShapira2008}, this holds for every hereditary graph
property: every such property is \emph{testable}, and in particular every
sequence with vanishing density of the forbidden configurations can be
$o(n^2)$-perturbed to a sequence in the constrained class.
\end{remark}

\section{The Case of General \texorpdfstring{$\sigma$}{sigma}}

For non-trivial $\sigma$, the direction $\ref{cond:quot} \Rightarrow
\ref{cond:rand}$ remains valid by the argument above. The reverse direction
is more delicate.

\begin{question}\label{q:gen}
For $\sigma \neq \emptyset$, does~\ref{cond:rand} imply~\ref{cond:quot}?
\end{question}

The obstruction is that not every $\psi \in \Homp(\A^{\sigma}[\T'], \R)$ need
arise (in any precise sense) as a member $\phi^{\sigma}$ of an ensemble
$\RandomExt(\phi_0; \sigma)$ rooted at some
$\phi_0 \in \Homp(\A^{\emptyset}[\T], \R)$ satisfying $A$. The natural candidate
for such a $\phi_0$ is $\Phi^{\sigma}(\psi \circ \pi^{\sigma})$ from
Section~\ref{sec:bridge}\,(b), and then Theorem~3.17 of~\cite{Razborov2007}
identifies $\RandomExt(\phi_0; \sigma)$ as the canonical disintegration of
$\phi_0$ along $\Phi^{\sigma}$. The question is whether the particular
homomorphism $\psi$ is in the support of this disintegration; in general it
need not be, and this gap is the source of the difficulty.

A reasonable conjecture, motivated by the discussion above, is:

\begin{conjecture}
The equivalence of Question~\ref{q:main} holds for arbitrary $\sigma$ if
\ref{cond:rand} is strengthened to:
\begin{enumerate}[label=\textup{(2$'$)}, leftmargin=3em]
  \item For every $\phi_0 \in \Homp(\A^{\emptyset}[\T], \R)$ satisfying $A$
        and $\phi_0(\langle \sigma \rangle) > 0$, and for every
        $\psi$ in the support of $\RandomExt(\phi_0; \sigma)$, we have
        $\psi(f) \geq 0$.
\end{enumerate}
This is in turn equivalent to requiring $\psi(f) \geq 0$ for every
$\psi \in \Homp(\A^{\sigma}[\T'], \R)$ that lies in the support of some
ensemble.
\end{conjecture}

A complete answer to Question~\ref{q:gen} would presumably require a
clean characterisation of which $\sigma$-homomorphisms of $\T'$ are
realisable as ensemble members; this is closely related to the open
questions raised in Razborov's~\S6.

\section*{Acknowledgements}

This note grew out of a series of conversations with an AI assistant about
Razborov's flag-algebra paper~\cite{Razborov2007}, in particular the role of
the induced removal lemma in interpreting quotient semantics.

\begin{thebibliography}{9}

\bibitem{Razborov2007}
A.~A.~Razborov,
\emph{Flag Algebras},
Journal of Symbolic Logic \textbf{72} (2007), no.~4, 1239--1282.

\bibitem{AFKS2000}
N.~Alon, E.~Fischer, M.~Krivelevich, and M.~Szegedy,
\emph{Efficient testing of large graphs},
Combinatorica \textbf{20} (2000), no.~4, 451--476.

\bibitem{AlonShapira2008}
N.~Alon and A.~Shapira,
\emph{A characterization of the (natural) graph properties testable with
one-sided error},
SIAM Journal on Computing \textbf{37} (2008), no.~6, 1703--1727.

\bibitem{LovaszSzegedy2006}
L.~Lov\'asz and B.~Szegedy,
\emph{Limits of dense graph sequences},
Journal of Combinatorial Theory, Series B \textbf{96} (2006), no.~6, 933--957.

\end{thebibliography}

\end{document}
